\hypertarget{group-state-and-general-measurement}{%
\section{Group state and general measurement}\label{method-insurance-methods-group-state-and-general-measurement}}

The following cases apply selected general-measurement and service-allocation provisions from the retained AASB 17 compilation \citep{aasb17}. They extend the premium-allocation example without presuming actuarial calibration or a complete insurer model.

The general method distinguishes fulfilment cash flows (FCF), including an
explicit risk adjustment (RA), from the contractual service margin (CSM),
which holds unearned profit. The liability for remaining coverage (LRC)
concerns future service; the liability for incurred claims (LIC) concerns
service and insured events already experienced. Keeping these components
separate makes a change in future estimates different from an unpaid claim.

Contract records retain insured risks, premiums/benefits, service and settlement dates, options/guarantees, cancellation/repricing rights and acquisition costs. Group admission evaluates portfolio/similar-risk management, onerous status and required issuance cohorts under AASB 17.14--24, recognition under 25--28 and boundary under 34. Cash outside the boundary belongs to future contracts, not this group's profit. GroupState retains expected dated cash flows, current discount curves, inception rates, explicit risk adjustment, CSM, loss component, LRC/LIC, coverage units, service experience and assumption vintages. Same premium totals are not sufficient state for merging different guarantees or claim tails.

General measurement follows 32--38: fulfilment cash flows are probability-weighted future net outflows, adjusted for timing/financial risk without double counting, plus explicit non-financial risk adjustment. CSM removes initial unearned gain under its conditions; an initial net loss is recognized under 47, not negative CSM. Paragraphs 40--44 separate remaining coverage and incurred claims. For a nonparticipating group, accrete CSM at the specified locked-in rate, apply future service changes at their required rate to CSM up to its available amount, record any excess loss and allocate current service over current plus future coverage units under B119. B96--B97 distinguish future service from incurred-claim changes, financial changes and current experience; not every revision is a CSM adjustment.

A MeasurementInput must resolve cash-flow/law, discount, risk-adjustment method, service grouping and coverage-unit evidence. Missing probabilities or RA support permit conditional sensitivity cases but block a source-supported measurement. Risk adjustment is the compensation the entity requires under 37; a standard deviation, confidence quantile or arbitrary percentage is only a candidate method until justified for that entity. Disclosures and the confidence-level equivalent, where applicable, are separate obligations. Market-consistent financial inputs and entity risk compensation must not be conflated.

\hypertarget{complete-two-year-general-measurement-case}{%
\section{Complete two-year general-measurement case}\label{method-insurance-methods-complete-two-year-general-measurement-case}}

Take a two-year nonparticipating group with one premium 120 received at inception, claims expected/actually incurred and paid 40 each year, zero discount rate, RA 10 initially released 5 with each equal service year, no acquisition costs, investment component, lapses, taxes, finance changes or reinsurance. Coverage units are 1 in each year because equal benefits and coverage are assumed. Before premium payment, expected net cash outflow is 80\ensuremath{-}120=\ensuremath{-}40, RA 10, FCF \ensuremath{-}30 and CSM 30, with net initial carrying amount zero. After the premium debit cash 120/ credit LRC 120. Its components are remaining expected claims 80, RA 10 and CSM 30.

Each year releases expected service claims 40, RA 5 and CSM 15: debit LRC 60/ credit insurance revenue 60. Incur claims 40: debit service expense 40/credit LIC 40; payment debits LIC 40/credits cash 40. End of year one: cash 80, LRC 60 (claims 40, RA 5, CSM 15), LIC 0 and equity/profit 20. End year two: cash 40, liabilities 0, cumulative equity/profit 40. Revenue 120 and claims 80 reconcile to premium 120 less claims cash 80. No investment component has been mistaken for revenue.

At year-one close, a separate pre-next-service change increases the remaining claim estimate from 40 to 60, RA unchanged. FCF rises 20; consume remaining CSM 15 and recognize service loss 5. LRC rises from 60 to 65, comprising FCF 65, CSM 0 and loss component 5 as a tracked portion, not an added liability on top. Debit service loss 5/credit LRC 5 with component transfer CSM\ensuremath{\rightarrow}FCF 15. A decrease of 8 in the revised future claim estimate, before any additional service/finance movement, first reverses loss 5 and restores CSM 3. FCF 57 plus CSM 3 gives LRC 60. Credit loss reversal 5/debit LRC 5. This tests 48--52 ordering. Later systematic allocation of loss component to releases/finance must reduce it to zero by coverage end; it cannot be amortized as revenue a second time.

In a separate case, an already-incurred unpaid claim estimate rises 40 to 48. Debit service expense 8/credit LIC 8; do not consume future-service CSM. A discount-rate change is financial under B97; its profit/OCI route follows the selected permitted finance policy. The zero-rate base case does not remove the need for separate current and locked-in rate calculations in a real group.

For the unreversed onerous alternative, carry the loss component through the remaining service year. Opening LRC 65 comprises future claims 60 and RA 5, with CSM 0 and loss component 5 within that LRC. A declared proportional allocation assigns 5/65 of each expected claim/RA release to the loss component, consistent with paragraphs 49--52 and B124's revenue exclusions. The expected claim release allocates 60/13 to loss reversal and the RA release 5/13; the total is 5, leaving loss component 0. Debit LRC 65, credit revenue 60 and credit loss reversal 5. Separately incur and pay claims 60 through LIC. Period profit is 60\ensuremath{-}60+5=5; combined with first-year profit 20 and the earlier onerous loss 5, lifetime profit is 20. Closing cash 20 equals equity 20 and all liabilities are zero. Treating all 65 as revenue while also reversing loss 5 overstates profit by 5.

\hypertarget{nonzero-discount-and-service-allocation}{%
\section{Nonzero discount and service allocation}\label{method-insurance-methods-nonzero-discount-and-service-allocation}}

A second complete case keeps premium 120, two equal coverage units and RA 10, but expected and realized claims are 44 after one year and 48.4 after two years. Both the inception and current discount curves remain a flat 10\% throughout; there are no new estimates, options, fees, taxes, investments or OCI election. This RA schedule is an authored measurement input, not an actuarial calibration: it declines by 5 per year, with its entire change in the service result under 81. All other insurance finance effects go to profit or loss. The retained source supports the mechanics in 44, 80--88, B119 and B123--B124.

At inception the present value of the two expected claims is
\[
 \frac{44}{1.1}+\frac{48.4}{1.1^2}=40+40=80.
\]
 FCF before premium is 80+10\ensuremath{-}120=\ensuremath{-}30, CSM 30, and after premium LRC 120. In year one the claim PV accretes 8; CSM accretes 3 at the locked-in rate. Allocate the resulting CSM 33 across the two current/future units, releasing 16.5. Revenue is expected claims 44 plus RA 5 plus CSM 16.5=65.5. Service expense 44 and finance expense 11 leave profit 10.5. After paying 44, cash 76 and LRC 65.5 (claim PV 44, RA 5, CSM 16.5) reconcile to equity 10.5.

In year two claim PV accretes 4.4 and CSM accretes 1.65. Release all remaining CSM 18.15, RA 5 and expected claims 48.4, giving revenue 71.55. Service expense 48.4 and finance expense 6.05 leave profit 17.1. Final cash/equity 27.6 equals premium 120 minus claims 92.4; all group liabilities are zero. Total revenue 137.05 less finance expense 17.05 equals premium 120. Setting revenue equal to cash premium each year or releasing 15 of CSM each year would ignore the specified financing. These are separate component calculations followed by debit finance expense/ credit LRC for accretion, debit LRC/credit revenue for service and the stated LIC recognition/settlement entries. Cash is deliberately assumed uninvested.

\hypertarget{paa-reinsurance-and-valuation-methods}{%
\section{PAA, reinsurance and valuation methods}\label{method-insurance-methods-paa-reinsurance-and-valuation-methods}}

The preceding premium-allocation case states separate eligibility and recognition timing for gross insurance and reinsurance. PAA under 53--59 requires the one-year coverage or nonmaterial-difference condition; a short payment lag is not the eligibility test. LRC starts with received premium and applicable acquisition adjustments, then releases insurance revenue under the prescribed service pattern. LIC still measures incurred obligations using the applicable fulfilment-cash-flow rules; the allowed short-settlement discount simplification needs its own conditions. Emerging onerous facts require a loss assessment; a premium balance cannot be presumed sufficient.

A one-year policy can leave claims unpaid for several years. The distinction
between the period of insurance service and the period until settlement is
therefore essential even when PAA is available. The official ASBE~25 text
separates three timing tests. Article~50 permits PAA when each contract's
coverage is no longer than one year, or when the stated no-material-difference
condition holds. Article~54 concerns a different interval: the time between
each portion of service and the related premium due date. If that interval
does not exceed a year, the remaining-coverage liability need not be adjusted
for financing. Article~56 instead measures from claim incurrence to expected
payment or receipt. Its option to omit the time-value and financial-risk
adjustment requires settlement within a year of incurrence
\citep[Articles 50, 54, 56]{mofasbe25}. The retained AASB~17 compilation makes
the corresponding distinctions in paragraphs 53, 56 and 59(b).

Consider a separate group of issued policies with exactly one year of coverage.
All amounts are illustrative currency units. Premium 120 is due and received
when coverage starts, with no earlier onerous trigger; group recognition is
therefore at that date under Article~16. Assume valid insurance grouping,
uniform risk release, no acquisition costs, investment component, reinsurance,
tax, other expenses or other movements, and no facts indicating onerous
remaining coverage. Elect the permitted remaining-coverage financing exception.
After half a year, recognise revenue 60 and retain remaining coverage of 60.
Articles~53 and 57 support this roll-forward and time allocation; Article~57
requires a different allocation when risk release is not predominantly
time-based \citep{mofasbe25}. These assumptions isolate the claim measurement
from other changes in the group's result.

At that reporting instant a claim is first incurred, with an expected payment
of 60.5 exactly two years later. This is a mean estimate, not a certain
settlement. Assume an appropriate flat annual effective discount rate of
10\% and a separately estimated non-financial risk adjustment of 5 in current
measurement units. The inputs are authored for the example; applying them to
an insurer would require market and actuarial evidence. Articles~21, 23, 25
and 26 require current probability-weighted cash flows, appropriate discounting
and a separate non-financial risk adjustment \citep{mofasbe25}. The claim's
cash-flow present value is $60.5/1.1^2=50$, so its liability is $50+5=55$.
The risk adjustment is not discounted a second time. The claim has only just
been incurred, so no intervening finance accrual enters this measurement.

Record debit cash 120/credit remaining coverage 120 on premium receipt, then
debit remaining coverage 60/credit insurance revenue 60 for the service
provided. Recognising the claim gives debit insurance service expense 55/credit
incurred claims 55, consistent with Article~31. The group then has cash 120,
remaining coverage 60, incurred claims 55 and a cumulative result of 5:
$120-60-55=60-55=5$. These are balances attributable to the group at mid-year,
not a complete insurer forecast. Using the undiscounted $60.5+5=65.5$ for this
two-year claim would overstate the required liability by 10.5 and understate
the result by the same amount. The shortcut is wrong under the stated rules
and assumptions. A zero rate would remove that discounting difference while
leaving the separate risk adjustment intact. Articles~86(3) and 97 also require
disclosure of the distinct claim components and the PAA/financing choices
\citep{mofasbe25}.

This comparison uses identified editions. The Ministry page reproduces the
2020 ASBE~25 text; Article~120 states 1 January 2023, but that sentence does
not settle every entity group's adoption schedule.
Section~\ref{sec:insurance-adoption} distinguishes the Ministry's later
explanation of that timetable from a company's reported adoption. The retained
December 2022 AASB~17 compilation applies to annual periods beginning on or
after 1 January 2023 but before 1 July 2026. The shared result in this example
establishes a scoped measurement comparison. Applying it to a particular entity
still requires its applicable rules and contract facts; later amendments and
broader equivalence require separate evidence.

Reinsurance has its own group, boundary, recognition time, counterparty risk, cash-flow, CSM/cost and possibly loss-recovery component under AASB~17 paragraphs 60--70A. Do not multiply every gross insurance liability by a quota share: excluded claims, limits, deductibles, timing, reinstatements and counterparty non-performance matter. Claim 100 with an eligible 50\% recovery but expected 90\% collectability gives conditional expected recovery 45, not 50; altering transferred insurance risk and reinsurer non-performance are different inputs. This simple calculation does not establish treaty recognition or RA by itself. Gross liability 100 remains; reinsurance recovery is a separate asset/measurement with its own authority.

For direct participation, test B101's conditions using the actual underlying items, substantial-share and variation judgments. Retain underlying fair values, entity fee share and guarantee/risk-mitigation choices; paragraphs 45 and B101--B118 specify the variable-fee adjustments to CSM, with onerous and mitigation exceptions. A local component illustration of a qualifying share: underlying-value increase 20 and entity's fixed 10\% share gives candidate variable-fee change 2; eligibility, guarantees and other FCF changes must be evaluated before posting. The full option/guarantee valuation and transition calculations are further supported profiles, with source and numerical obligations, not an assumed linear fee model.

Eligibility depends on the contract and the insurer's expectations at inception.
The policyholder must participate in a clearly identified pool, with the
required substantial share of fair-value returns and substantial variability
in payments. The variability assessment covers the contract's duration on a
probability-weighted present-value basis. A minimum guarantee can bind in some
scenarios without automatically excluding direct participation. Later market
movements do not reopen the classification; contract modification follows its
own rules. The identified reference pool also need not be held by the insurer,
so the contractual link and the actual asset holdings require separate records
\citep[paragraphs B101--B108 and 72]{aasb17}.

\clearpage
A small settlement example shows why the guarantee matters even when the
stated share is fixed. Consider only a maturity-benefit component, in one
currency at one date. Assume a held pool allocated to this component has amount
$U$, and the policyholder is owed the larger of 90 and $0.9U$. These are
illustrative contractual choices. The amount left from the pool after payment
is $U-\max(90,0.9U)$. Writing the payment as
$0.9U+\max(90-0.9U,0)$ shows that the amount left is
$0.1U-\max(90-0.9U,0)$: the nominal retained share less the extra payment
needed to meet the minimum.

\begin{center}
\begin{tabular}{@{}rrr@{}}
\toprule
Pool amount & Policyholder payment & Amount left\\
\midrule
80 & 90 & $-10$\\
100 & 90 & 10\\
120 & 108 & 12\\
\bottomrule
\end{tabular}
\end{center}

Moving from 80 to 100 leaves 20 more after payment: the ordinary share rises
by 2 and the extra floor payment falls by 18. Moving from 100 to 120 leaves
only 2 more because the minimum no longer binds. The negative amount at 80
is a funding shortfall for this component. These unweighted settlement
outcomes explain the exposure; they do not establish insurance-contract
eligibility, a present value or the insurer's profit. Other benefits and costs
remain outside this component example, and its held-pool assumption does not
replace the standard's separate accounting for investment assets
\citep[paragraphs B66(a), B107--B108]{aasb17}.

The accounting calculation therefore needs more than this payoff schedule.
Paragraph B104 defines the variable fee as the entity's share of underlying
fair value less fulfilment cash flows that do not vary based on underlying
returns. Changes in the obligation to pass through the underlying fair value
do not adjust CSM, whereas the entity-share changes relate to future service.
Financial-guarantee changes have their own future-service treatment under
B113(b). Paragraph 45 retains the loss-component and risk-mitigation exceptions,
as well as new-contract, currency and service-allocation movements. Applying
those rules requires the group's measured cash flows, discount and risk
adjustments, remaining CSM and loss component, service information and any
eligible mitigation designation. The original increase of 2 remains a valid
share-component calculation; the complete CSM movement needs these additional
inputs \citep[paragraphs 32--38, 45, 48--50 and B104--B118]{aasb17}.

The claims-development example above supplies an estimation input, not a
recognition rule. Its point projection is 148, but the single observed final
factor leaves its variance unresolved. In the general method the factors are
\[
 \widehat f_k=\frac{\sum_{i\in O_k} C_{i,k+1}}{\sum_{i\in O_k}C_{i,k}},
\]
where $O_k$ contains origin years with both required observed cells. Each
incomplete row is projected through the remaining factors. The same triangle
therefore has ultimate claims 180, 216 and 252, with unpaid amounts 0, 36 and
112. Tail development, calendar inflation, zero or negative cells and dependence
need explicit extensions. The conditional means, variance scaling and
independent-origin-year assumptions behind the inspected Mack error formulas
\citep[equations (1)--(3), Theorems 1--3]{mack1993} cannot turn this tiny
triangle into a predictive law or a justified insurance risk adjustment.

\section{Adoption dates and the restatement of 2022}
\label{sec:insurance-adoption}

A standard's effective date and the economic date of a reported balance need
not coincide. An insurer adopting new rules in 2023 may report a revised
balance for December 2022 while retaining the earlier report as a record of
what it originally disclosed. The system must preserve both versions and
explain which one a calculation uses.

The Ministry of Finance's explanation dated 12 June 2025 describes the original
ASBE~25 timetable. It specifies application from 1 January 2023 for enterprises
listed both domestically and overseas, and for overseas-listed enterprises
preparing financial statements under IFRS or ASBE. Other enterprises applying
ASBE were scheduled to apply the revised standard from 1 January 2026.
Its account of the
2025 notice says existing adopters, including enterprises that had adopted
early under the 2020 notice, must continue applying it. Insurers facing genuine
difficulties and seeking deferral must submit written reasons to the Ministry's
Accounting Department and the relevant institutional supervision departments
of the National Financial Regulatory Administration. That statement does not
establish an automatic deferral merely because preparation is difficult
\citep[answers on background and implementation arrangements]{mofinsurance2025qa}.

The same explanation limits the new simplification choices to unlisted
enterprises; matters outside those choices remain subject to ASBE~25. It names
grouping, measurement, presentation and transition as affected areas, including
comparative information and the transition date. These are source-supported
descriptions of the Ministry's explanation. The operative 2020 and 2025 notices
have not yet been directly verified in this collection, so the explanation
does not supply the detailed conditions needed to implement those choices
\citep[answer on the relationship between the notice and the standard]{mofinsurance2025qa}.

Ping An provides a concrete adoption history. Its 2022 ASBE annual report
states that the group would adopt the revised insurance standard for accounting
years beginning on 1 January 2023. Its 2023 report confirms adoption on that
date and restatement of comparative information
\citep[Note XIV.2, pp.~315--316]{pingan2022asbe}
\citep[Note IV.34, p.~192]{pingan2023asbe}. The following balances all refer
to the consolidated group at 31 December 2022, in RMB millions. The first
column comes from the 2022 report; the second is the restated comparative in
the 2023 report. Changes are calculated from those columns; the asset,
liability and parent-equity changes also match the three rows of Note IV.34's
reported bridge \citep[pp.~146--148]{pingan2022asbe}
\citep[pp.~141--143 and 192]{pingan2023asbe}.

\begin{center}
\small
\begin{tabularx}{\linewidth}{@{}Xrrr@{}}
\toprule
Balance & Original & Restated & Change\\
\midrule
Assets & 11,137,168 & 11,009,940 & $-127,228$\\
Liabilities & 9,961,870 & 9,823,944 & $-137,926$\\
Total equity & 1,175,298 & 1,185,996 & $10,698$\\
Parent shareholders' equity & 858,675 & 869,191 & $10,516$\\
Non-controlling interests & 316,623 & 316,805 & $182$\\
\bottomrule
\end{tabularx}
\end{center}

The distinction between total equity and equity attributable to the parent
matters. Assets less liabilities rise by
$-127,228-(-137,926)=10,698$. Of that increase, 10,516 belongs to parent
shareholders and 182 to non-controlling interests. Treating the parent-equity
row as total equity would leave an unexplained difference of 182. These are
changes between reported accounting bases for the same date, not cash received
in a later period.

The transition also changes how later growth is calculated. Ping An reports
assets of 11,583,417 at the end of 2023. Subtracting its restated 2022
comparative gives a reported-balance increase of
$11,583,417-11,009,940=573,477$. Subtracting the original 2022 amount instead
gives 446,249, mixing the old and new bases and understating that comparable
movement by 127,228 \citep[p.~141]{pingan2023asbe}. The first difference is
a movement in reported asset balances; it is not a cash-flow measure or an
explanation of every economic cause of growth.

An original-source query for assets at 31 December 2022, explicitly selecting
the 2022 report, therefore returns 11,137,168. A query selecting the 2023
report's restated comparison for that date returns 11,009,940. Each assertion
retains its report identity, consolidated scope, RMB-million scale, accounting
basis and economic date. A link records that the later assertion restates the
comparison; it does not delete the original. A request for a comparable annual
movement must reject the mixed-basis pair unless an explicit reconciliation
has first aligned it. Because the basis and disclosure version differ, these
two amounts are not a same-context source conflict.

The information date remains a separate question. The 2023 report says its
board authorised the statements on 21 March 2024
\citep[Note I, p.~154]{pingan2023asbe}. Neither that authorisation nor the
31 December reporting date establishes when the exact retained document became
public. The corpus has archive-date leads but no verified publication date or
historical byte vintage for either selected report. A query for information
available to investors at 31 December 2023 must therefore retain an unresolved
availability result for this source selection; it cannot silently import the
later restated balance. The corpus's October 2026 revalidation timestamps
describe local evidence handling, not investor knowledge in 2023.

Finally, Note IV.34 reports financial-asset reclassification, including use of
the classification overlay for relevant assets derecognised during 2022. It
also states that some insurance groups used modified-retrospective or
fair-value transition because retrospective application was impracticable.
The headline bridge consequently cannot be reconstructed by applying one
percentage to every insurance liability. The example establishes how to retain
and compare disclosed histories; reconstructing the transition itself still
requires the affected groups, asset classifications and underlying measurement
inputs \citep[Note IV.34, p.~192]{pingan2023asbe}.

\clearpage
\hypertarget{interfaces-exclusions-and-acceptance-cases}{%
\section{Interfaces, exclusions and acceptance cases}\label{method-insurance-methods-interfaces-exclusions-and-acceptance-cases}}

Insurance events are group\_recognize, premium\_settle, service\_release, claim\_incur, claim\_settle, assumption\_change, onerous\_change and treaty\_change, each bound through the event-basis, authorization and commit contracts to the exact group/version, evidence and information view. Journals and component roll-forwards must reconcile independently; a component reallocation is not always a net liability change. Invalid grouping, missing rate/law, unsupported option valuation, ambiguous contract boundary or unknown service units produces the corresponding unresolved/outside-scope result. Failed numerical paths retain their coverage/mass, without successful-only normalization. Modified contracts need the applicable derecognition analysis; no new ID can erase unfavorable old-group obligations.

Initial transition requires an actual method choice and eligibility for full retrospective, modified-retrospective or fair-value transition under AASB~17 Appendix C, with acquisition, grouping, past service and historical-rate records. This collection does not claim an executed transition or a general VFA valuation. The nonzero-rate general case and onerous allocation through expiry above are finite specification examples. Required later acceptance still includes a priced guarantee, acquisition/reinsurance interaction, transition and actuarial backtesting. These retained obligations prevent the two-year illustration from being described as full insurance-domain readiness.
